\paragraph{Fixed regular baseline.}
Let
\[
\mathcal E=\{-1,+1\},\qquad
\mathcal U=\mathcal W=\mathcal Y=\{0,1\},\qquad
\mathcal X=\{x\}.
\]
Order vectors on each binary state space as \((1,0)\), the environment columns as \((+1,-1)\), and the latent columns of the proxy channel as \((1,0)\). Define \(\mathcal M^{(0)}\) by
\[
\pi=(1/2,1/2),\qquad
S^{(0)}=
\begin{pmatrix}1/2&1/2\\1/2&1/2\end{pmatrix},\qquad
M^{(0)}=
\begin{pmatrix}3/4&1/4\\1/4&3/4\end{pmatrix},
\]
\[
a_x(u)=1,\qquad q^{(0)}=(1/2,1/2)^{\top},\qquad
f^{(0)}_{x,1}(u,w)=f^{(0)}_{x,0}(u,w)=1/2.
\]
Every primitive entry is strictly positive, every displayed kernel is stochastic, and
\[
\det M^{(0)}=\frac12\ne0.
\]
Thus the two factorization assumptions, strict primitive positivity, proxy-channel injectivity, and the definition of \(\mathfrak M_{+}\) give
\(\mathcal M^{(0)}\in\mathfrak M_{+}\).

Summing the source factorization over \(u\) gives
\[
P_{O,0}(e,w,x,y)
=\sum_u\pi_eS^{(0)}_{ue}M^{(0)}_{wu}a_x(u)f^{(0)}_{x,y}(u,w)
=\frac18
\]
for every \(e,w,y\). Target-mechanism invariance gives
\[
b_0=M^{(0)}q^{(0)}=(1/2,1/2)^{\top}
\]
and
\[
\theta_{x,1}(\mathcal M^{(0)})
=\sum_{u,w}q^{(0)}_uM^{(0)}_{wu}f^{(0)}_{x,1}(u,w)
=\frac12.
\]
By the definition of the unconditional moments,
\[
H_{x,0}=
\begin{pmatrix}1/4&1/4\\1/4&1/4\end{pmatrix},
\qquad
z_{x,1,0}=(1/4,1/4)^{\top}.
\]
The singular values of \(H_{x,0}\) are \(1/2\) and \(0\), so
\[
\operatorname{rank}(H_{x,0})=1,\qquad
\sigma_1(H_{x,0})=\frac12\ge\frac14,\qquad
b_0=H_{x,0}(1,1)^{\top}.
\]
Every source observed cell is \(1/8\), and every target proxy cell is \(1/2\); hence the cell floor is at least \(1/16\).

We verify the remaining positive-variance condition. Since \(H_{x,0}\) has rank one,
\[
H_{x,0}^{\dagger}=
\begin{pmatrix}1&1\\1&1\end{pmatrix},
\qquad
D_z\phi_1(H_{x,0},z_{x,1,0},b_0)
=\bigl(H_{x,0}^{\dagger}b_0\bigr)^{\top}
=(1,1).
\]
The source part of the scalar delta-method influence variable therefore has the form
\(\mathbf 1\{Y=1\}+\psi(E,W)\), up to an additive constant, for a deterministic function \(\psi\) contributed by the \(H\)-coordinates. Under \(P_{O,0}\), the constant outcome kernel makes \(Y\) conditionally fair given \((E,W)\). Expanding the finite conditional sums gives the variance decomposition
\[
\begin{split}
\operatorname{Var}_{P_{O,0}}\!\left(\mathbf 1\{Y=1\}+\psi(E,W)\right)
&=\mathbb E_{P_{O,0}}\!\left[
\operatorname{Var}_{P_{O,0}}\!\left(\mathbf 1\{Y=1\}\mid E,W\right)
\right]\\
&\quad+
\operatorname{Var}_{P_{O,0}}\!\left(
\mathbb E_{P_{O,0}}[
\mathbf 1\{Y=1\}+\psi(E,W)\mid E,W]
\right)\\
&\ge\frac14.
\end{split}
\]
The source allocation factor is \(p^{-1}\), and the independent target contribution is nonnegative. Consequently,
\[
V_{1,x,1}(\mathcal M^{(0)};p)\ge\frac1{4p}>\frac1{16}.
\]
All five defining conditions of the strongly identified submodel have been checked, and hence
\[
\mathcal M^{(0)}\in\mathcal R_{1,1/4,1/16}.
\]

\paragraph{Admissible weak-rank array.}
For \(n\ge2\), set \(\gamma_n=\delta_n=n^{-2}\). Define \(\mathcal M^{(1)}_n\) using the same \(\pi\) and \(a_x\), and put
\[
S^{(1)}_n=
\begin{pmatrix}
1/2+\gamma_n&1/2-\gamma_n\\
1/2-\gamma_n&1/2+\gamma_n
\end{pmatrix},\qquad
M^{(1)}_n=
\begin{pmatrix}
1/2+\delta_n&1/2-\delta_n\\
1/2-\delta_n&1/2+\delta_n
\end{pmatrix},
\]
\[
q^{(1)}=(3/4,1/4)^{\top},\qquad
f^{(1)}_{x,1}(1,w)=3/4,\qquad
f^{(1)}_{x,1}(0,w)=1/4,\qquad
f^{(1)}_{x,0}=1-f^{(1)}_{x,1}.
\]
For \(n\ge2\), all entries of \(S^{(1)}_n\), \(M^{(1)}_n\), \(q^{(1)}\), and \(f^{(1)}\) belong to \([1/4,3/4]\), while \(a_x=1\) and \(\pi_e=1/2\). Moreover,
\[
\det M^{(1)}_n=2\delta_n>0,\qquad
\det S^{(1)}_n=2\gamma_n>0.
\]
Therefore every row satisfies strict primitive positivity and proxy-channel injectivity, and the factorization assumptions and the definition of \(\mathfrak M_{+}\) give
\(\mathcal M^{(1)}_n\in\mathfrak M_{+}\).

Because \(x\) is the sole treatment value, summing the source factorization yields
\[
B^{(1)}_{x,n}=M^{(1)}_nS^{(1)}_n=
\begin{pmatrix}
1/2+2\delta_n\gamma_n&1/2-2\delta_n\gamma_n\\
1/2-2\delta_n\gamma_n&1/2+2\delta_n\gamma_n
\end{pmatrix},
\]
\[
H^{(1)}_{x,n}=\frac12B^{(1)}_{x,n},
\qquad
\det B^{(1)}_{x,n}=4\delta_n\gamma_n>0.
\]
Target-mechanism invariance gives
\[
b^{(1)}_n=M^{(1)}_nq^{(1)}
=
\begin{pmatrix}1/2+\delta_n/2\\1/2-\delta_n/2\end{pmatrix}.
\]
Since \(B^{(1)}_{x,n}\) is invertible,
\(b^{(1)}_n\in\operatorname{col}(B^{(1)}_{x,n})\) for every \(n\). Together with the triangular allocation and row-wise independent sampling assumptions, this verifies
\[
\mathbb M^{(1)}=(\mathcal M^{(1)}_n)_{n\ge2}\in\mathcal A_p.
\]
The array reaches the unrestricted weak-rank part of the class. If \(\lambda_n\) is the unique solution of
\(B^{(1)}_{x,n}\lambda_n=b^{(1)}_n\), addition and subtraction of the two row equations give
\[
\lambda_{1,n}+\lambda_{2,n}=1,\qquad
\lambda_{1,n}-\lambda_{2,n}=\frac1{4\gamma_n}.
\]
Thus \(\lVert\lambda_n\rVert_2\to\infty\). Since
\(P_{O,n}(E=e,X=x)=1/2\), the corresponding unconditional balancing-weight norm also diverges.

The causal value remains separated from the baseline. The target factorization and the intervention definition give
\[
\begin{split}
\theta_{x,1,n}(\mathcal M^{(1)}_n)
&=\sum_{u,w}q^{(1)}_uM^{(1)}_n(w,u)f^{(1)}_{x,1}(u,w)\\
&=\frac34\cdot\frac34+\frac14\cdot\frac14
=\frac58
\end{split}
\]
for every \(n\).

\paragraph{Indistinguishability of the experiments.}
For probability vectors \(P,Q\) on a common finite space, write
\[
\lVert P-Q\rVert_{\mathrm{TV}}=\frac12\sum_a|P(a)-Q(a)|.
\]
As an auxiliary comparison, retain \(f^{(1)}\), but replace both
\(S^{(1)}_n\) and \(M^{(1)}_n\) by the binary all-\(1/2\) kernels
\(S^\star\) and \(M^\star\). The resulting observed law of
\((E,W,X,Y)\) equals \(P_{O,0}\): \(E\) and \(W\) are independent and fair, and averaging the response probabilities \(3/4\) and \(1/4\) over the fair latent variable makes \(Y\) fair independently of \((E,W)\).

At the full-data level, add and subtract the density with
\(S^\star\) and \(M^{(1)}_n\):
\[
\begin{split}
&\pi_eS^{(1)}_n(u,e)M^{(1)}_n(w,u)f^{(1)}_{x,y}(u,w)
-\pi_eS^\star(u,e)M^\star(w,u)f^{(1)}_{x,y}(u,w)\\
&=\pi_e\bigl(S^{(1)}_n(u,e)-S^\star(u,e)\bigr)
M^{(1)}_n(w,u)f^{(1)}_{x,y}(u,w)\\
&\quad+\pi_eS^\star(u,e)
\bigl(M^{(1)}_n(w,u)-M^\star(w,u)\bigr)
f^{(1)}_{x,y}(u,w).
\end{split}
\]
After summing absolute values, stochasticity of the unchanged kernels shows that the full-data total variation is at most the sum of the two conditional binary-kernel total variations, namely \(\gamma_n+\delta_n\). Marginalization only groups summands before the absolute value and therefore cannot increase this bound. Hence
\[
\lVert P^{(1)}_{O,n}-P_{O,0}\rVert_{\mathrm{TV}}
\le\gamma_n+\delta_n=2n^{-2}.
\]
The displayed target proxy vectors give
\[
\lVert b^{(1)}_n-b_0\rVert_{\mathrm{TV}}=\delta_n/2.
\]

For arbitrary finite probability vectors, the telescoping identity
\[
\bigotimes_{i=1}^N P_i-\bigotimes_{i=1}^N Q_i
=
\sum_{j=1}^N
\left(\bigotimes_{i<j}P_i\right)\otimes(P_j-Q_j)\otimes
\left(\bigotimes_{i>j}Q_i\right)
\]
and the triangle inequality imply that product total variation is bounded by the sum of the coordinatewise total variations, because every probability vector has \(\ell_1\)-norm one. Since
\(n_s(n)+n_t(n)=n\),
\[
\begin{split}
&\left\lVert
(P^{(1)}_{O,n})^{\otimes n_s(n)}
\otimes(b^{(1)}_n)^{\otimes n_t(n)}
-
P_{O,0}^{\otimes n_s(n)}
\otimes b_0^{\otimes n_t(n)}
\right\rVert_{\mathrm{TV}}\\
&\qquad\le
n_s(n)(\gamma_n+\delta_n)+n_t(n)\delta_n/2
\le2n^{-1}\longrightarrow0.
\end{split}
\]
This proves total-variation convergence of the induced two-sample laws.

\paragraph{Transferred coverage.}
Let
\[
A_n=\{5/8\in\mathcal C_n\}.
\]
The assumed measurability makes \(A_n\) an event. Uniform coverage and
\(\mathbb M^{(1)}\in\mathcal A_p\) imply
\[
\liminf_{n\to\infty}
\mathbb P_{\mathcal M^{(1)}_n}^{\,n_s(n),n_t(n)}(A_n)
\ge1-\alpha.
\]
For every event, the difference of its probabilities under two finite laws is bounded by their total variation distance, as follows directly by splitting the defining absolute sum over the event and its complement. The preceding convergence therefore yields
\[
\liminf_{n\to\infty}
\mathbb P_{\mathcal M^{(0)}}^{\,n_s(n),n_t(n)}(A_n)
\ge1-\alpha. \tag{1}
\]

\paragraph{Localization of the fixed-model Wald interval.}
Under \(\mathcal M^{(0)}\), every coordinate of
\(\widehat H_{x,n}\), \(\widehat z_{x,1,n}\), and \(\widehat b_n\) is an average of independent indicator variables. For an average \(\widehat\mu\) of \(N\) indicator variables,
\[
\mathbb E[(\widehat\mu-\mu)^2]\le\frac1{4N},
\qquad
\Pr\{|\widehat\mu-\mu|>t\}
\le\frac1{4Nt^2},
\]
where the probability inequality follows from
\(\mathbf 1\{|\widehat\mu-\mu|>t\}\le
(\widehat\mu-\mu)^2/t^2\).
There are finitely many coordinates, and the allocation assumption gives
\(n_s(n)/n\to p\) and \(n_t(n)/n\to1-p\), both strictly positive. A finite union bound gives
\[
(\widehat H_{x,n},\widehat z_{x,1,n},\widehat b_n)
-(H_{x,0},z_{x,1,0},b_0)=O_P(n^{-1/2}). \tag{2}
\]

At \(H_{x,0}\), the eigenvalues of
\(H_{x,0}H_{x,0}^{\top}\) are \(1/4\) and \(0\). In a sufficiently small neighborhood of \(H_{x,0}\), the roots of the quadratic characteristic polynomial remain separated. On this neighborhood the leading projector is
\[
P_1(H)=
\frac{HH^{\top}-\lambda_2(H)I}
{\lambda_1(H)-\lambda_2(H)}.
\]
Thus the roots, \(P_1(H)\), the rank-one truncation \(P_1(H)H\), its pseudoinverse, \(\phi_1\), and \(D\phi_1\) are continuously differentiable with bounded derivatives after shrinking the neighborhood if necessary. The plug-in multinomial covariance is a polynomial function of empirical cell frequencies and is bounded on the finite probability simplex. Equation \((2)\), the definitions of the Wald functional, estimator, and totalized plug-in variance, and the finite-dimensional mean-value inequality now give
\[
\widehat\theta^{\,W}_{x,1,n}
=\frac12+O_P(n^{-1/2}),
\qquad
\widehat V_n=O_P(1). \tag{3}
\]
Because \(0<\alpha<1\), the standard-normal quantile is finite, so the untruncated Wald radius is \(O_P(n^{-1/2})\). It follows that
\(I^W_{n,1-\alpha}\) is nonempty with probability tending to one. Define on every sample
\[
\Delta_n=
\begin{cases}
\displaystyle
\sup_{\zeta\in I^W_{n,1-\alpha}}|\zeta-1/2|,
&I^W_{n,1-\alpha}\ne\varnothing,\\
1,&I^W_{n,1-\alpha}=\varnothing.
\end{cases}
\]
Clipping an interval to \([0,1]\) cannot increase its maximum distance from \(1/2\). Hence \((3)\) yields
\[
\Delta_n=O_P(n^{-1/2})=o_P(1). \tag{4}
\]

\paragraph{Quantitative impossibility.}
Fix \(\varepsilon\in(0,1/8)\) and let
\(G_n=\{\Delta_n\le\varepsilon\}\). Equation \((4)\) gives
\[
\mathbb P_{\mathcal M^{(0)}}^{\,n_s(n),n_t(n)}(G_n^c)
\longrightarrow0. \tag{5}
\]
On \(A_n\cap G_n\), the definition of \(\Delta_n\) implies that the Wald interval is nonempty. Since \(\mathcal C_n\) is nonempty and contains \(5/8\) on \(A_n\), the first directed term in the definition of Hausdorff distance gives
\[
\begin{split}
d_H(\mathcal C_n,I^W_{n,1-\alpha})
&\ge
\inf_{\zeta\in I^W_{n,1-\alpha}}|5/8-\zeta|\\
&\ge
\frac18-
\sup_{\zeta\in I^W_{n,1-\alpha}}|\zeta-1/2|\\
&=\frac18-\Delta_n
\ge\frac18-\varepsilon.
\end{split}\tag{6}
\]
Therefore, under the fixed baseline two-sample law,
\[
\begin{split}
&\mathbb P_{\mathcal M^{(0)}}^{\,n_s(n),n_t(n)}
\left\{I^W_{n,1-\alpha}\ne\varnothing,\ 
d_H(\mathcal C_n,I^W_{n,1-\alpha})
\ge\frac18-\varepsilon\right\}\\
&\qquad\ge
\mathbb P_{\mathcal M^{(0)}}^{\,n_s(n),n_t(n)}(A_n\cap G_n)\\
&\qquad\ge
\mathbb P_{\mathcal M^{(0)}}^{\,n_s(n),n_t(n)}(A_n)
-
\mathbb P_{\mathcal M^{(0)}}^{\,n_s(n),n_t(n)}(G_n^c).
\end{split}
\]
Taking lower limits and using \((1)\) and \((5)\) proves
\[
\liminf_{n\to\infty}
\mathbb P_{\mathcal M^{(0)}}^{\,n_s(n),n_t(n)}
\left\{I^W_{n,1-\alpha}\ne\varnothing,\ 
d_H(\mathcal C_n,I^W_{n,1-\alpha})
\ge\frac18-\varepsilon\right\}
\ge1-\alpha. \tag{7}
\]
Because the Wald interval is empty with probability tending to zero, the same lower bound holds after any totalization on empty-interval samples. Since \(1/8-\varepsilon>0\) and \(1-\alpha>0\), equation \((7)\) rules out convergence of the Hausdorff distance to zero in probability and therefore also rules out convergence after multiplication by \(\sqrt n\).

Finally, unrestricted uniform coverage by itself remains available. For each row of every
\(\mathbb M\in\mathcal A_p\), the target-span assumption gives
\(b_n\in\operatorname{col}(B_{x,n})\). Apply the already-established finite-sample projection-coverage theorem to \(\mathcal M_n\), \(n_s(n)\), and \(n_t(n)\). The row-indexed concentration projection set then has coverage at least \(1-\alpha\) in every row. Taking the infimum over \(\mathcal A_p\) and then the lower limit preserves this bound, including along changing-rank arrays.